\documentclass[10pt,a4paper]{amsart}
\usepackage[margin=19mm]{geometry}
\usepackage{amsmath,amssymb,mathtools}
\usepackage[scr=rsfs,cal=euler]{mathalfa}
\usepackage{xfrac}
\usepackage[unicode,hidelinks,bookmarks=false]{hyperref}
\newcommand{\ie}{i.e.\ }
\newcommand{\cf}{cf.\ }
\newcommand{\resp}{respectively\ }
\newcommand{\Subsection}{\subsection}
\providecommand{\nobreakdash}{\nobreak\hbox{-}}
\setlength{\emergencystretch}{2em}
\multlinegap0pt
\usepackage{mathtools}
\usepackage{amssymb}
\usepackage[shortlabels]{enumitem}
\usepackage{float}
\usepackage{tikz}
\usepackage{xcolor}
\newtheorem{definition}{Definition}
\newtheorem{lemma}{Lemma}
\title{Packing and Covering Cycles Through Prescribed Vertices}
\author{Hanzhi Bai, Jin Yan}
\date{Source: arXiv:2608.26772v1 (CC BY 4.0)}
\begin{document}
\maketitle

\noindent\emph{Source and licence.} Packing and Covering Cycles Through Prescribed Vertices. Version arXiv:2608.26772v1 (CC BY 4.0), distributed by arXiv under the Creative Commons Attribution 4.0 International licence, http://creativecommons.org/licenses/by/4.0/, as recorded in the arXiv licence metadata for this exact version and shown on the abstract page https://arxiv.org/abs/2608.26772v1. Full licence notice: \texttt{licenses/arxiv-2608.26772-CC-BY-4.0.txt}.\par\smallskip

\noindent\textit{Editorial scope.} Counted unit: the article's lemma lem:cycle-correspondence (source0.tex lines 395--403). The article's complete original proof of it is delivered with it (source0.tex lines 405--457, one \texttt{proof} environment of 53 lines). No mathematics is written, repaired or re-derived: the statement and the proof are the article's own lines, in the article's own notation, and no retained line is substituted or re-flowed.  Retained uncounted as the source-local context the proof needs: lines 236--245; lines 363--370.  Nothing was omitted from any retained span, so the package carries no omission record.  No retained line contains a citation, so the wrapper carries no bibliography and no bibliography entry.  No retained line contains a figure environment, an image-inclusion command or drawing code, so no image is delivered and the package declares no asset dependency.  Editorial formatting only, in addition: an AMS \texttt{amsart} preamble that restates the article's own declarations and macros used by the retained lines, namely the environments \texttt{definition}, \texttt{lemma} on separate counters, together with the standard packages those lines need.  The retained environments print in this document as Definition 1, Lemma 1 in order of appearance, whereas the article prints the same items with the numbering of its own full document.  The complete native source of the article as distributed in the arXiv e-print for this exact version is bundled unchanged as \texttt{sources/208659/source0.tex}.
\begin{definition}\label{def:bidirected}
A \emph{bidirected graph} $B=(H,\sigma)$ consists of a loopless multigraph
$H$, called its \emph{underlying multigraph}, and a signing $\sigma$ that
assigns a sign $\sigma(x,f)\in\{+,-\}$ to every \emph{half-edge} $(x,f)$,
where $(x,f)\in I(H)$.  A cycle of $H$ is a \emph{bidirected cycle} if, at
each of its vertices, the two half-edges belonging to the cycle have opposite
signs; two parallel edges may form a cycle of length two.  The terms cycle packing, $S$-cycle,
$S$-cycle transversal, $\nu_S$, and $\tau_S$ then refer only to bidirected
cycles.
\end{definition}

\begin{lemma}\label{lem:block-traversal}
Let $Q$ be a bidirected cycle in $B$.  If $Q$ contains $z^i$ for some
$z\in V(G)\cup E(G)$ and $i\in\{0,1\}$, then $Q$ contains the subpath
$pz^iz^{1-i}q$, where $p$ and $q$ are distinct incidence vertices.  Moreover,
at every
incidence vertex $p_{v,e}$ on $Q$, one incident edge of $Q$ leads to
$v^0$ or $v^1$ and the other leads to $e^0$ or $e^1$.
\end{lemma}

\begin{lemma}\label{lem:cycle-correspondence}
For a bidirected cycle $Q$ of $B$, let
$V_Q:=\{v\in V(G):\{v^0,v^1\}\subseteq V(Q)\}$ and
$E_Q:=\{e\in E(G):\{e^0,e^1\}\subseteq V(Q)\}$.  The subgraph of $G$ with
vertex set $V_Q$ and edge set $E_Q$ is a cycle, denoted by $\pi(Q)$.
Conversely, for every cycle $C$ of $G$, there is a bidirected cycle $Q$ of
$B$ satisfying $\pi(Q)=C$.  Moreover,
$|T\cap V(Q)|=|S\cap V(\pi(Q))|$.
\end{lemma}

\begin{proof}
Every edge of $B$ has an endvertex of the form $z^i$.  Hence every cycle of $B$
contains an object vertex.  By Lemma~\ref{lem:block-traversal}, whenever $Q$
contains one vertex in a pair associated with $z$, it contains both vertices
of that pair and the internal edge $b_z$.  At every incidence vertex, the same
lemma forces $Q$ to pass between a vertex pair and an edge pair.  Consequently,
the pairs encountered cyclically along $Q$ alternate between pairs associated
with vertices of $G$ and pairs associated with edges of $G$.

Let $e=uv\in E(G)$.  The pair $\{e^0,e^1\}$ is adjacent only to
$p_{u,e}$ and $p_{v,e}$.  If $\{e^0,e^1\}\subseteq V(Q)$, then
Lemma~\ref{lem:block-traversal} gives a subpath
$pe^ie^{1-i}q$ of $Q$ in which $p$ and $q$ are distinct incidence
vertices.  It follows that $\{p,q\}=\{p_{u,e},p_{v,e}\}$.  Applying the final
assertion of Lemma~\ref{lem:block-traversal} at $p_{u,e}$ and $p_{v,e}$ shows
that both vertex pairs associated with $u$ and $v$ occur on $Q$.

Similarly, if the vertex pair associated with $v$ occurs on $Q$, then
Lemma~\ref{lem:block-traversal} gives a subpath
$p_{v,e}v^iv^{1-i}p_{v,f}$ for two distinct edges $e$ and $f$ incident
with $v$.  Thus exactly two members of $E_Q$ are incident with each member of
$V_Q$.  Reading the associated pairs cyclically along $Q$ also shows that the
subgraph $(V_Q,E_Q)$ is connected.  Hence it is a connected $2$-regular
subgraph of $G$.  Since $G$ has neither loops nor parallel edges, it has at
least three vertices and is therefore a cycle.  The sets $V_Q$ and $E_Q$
determine it uniquely; this proves that $\pi(Q)$ is well defined.

For the converse, write $C=v_1e_1v_2e_2\cdots v_ke_kv_1$, where $k\geq3$,
$e_j=v_jv_{j+1}$, and all indices are taken modulo $k$; in particular,
$e_0=e_k$.  For each $j\in\{1,\ldots,k\}$, choose
$\alpha_j,\beta_j\in\{0,1\}$.  In cyclic order of $j$, concatenate the $2k$
paths $p_{v_j,e_{j-1}}v_j^{\alpha_j}v_j^{1-\alpha_j}p_{v_j,e_j}$ and
$p_{v_j,e_j}e_j^{\beta_j}e_j^{1-\beta_j}p_{v_{j+1},e_j}$.  Cyclically
consecutive paths have exactly their indicated common endvertex, while all
other vertices are distinct because $C$ is a cycle.  Their union is connected
and every vertex in it has degree two; it is therefore a cycle in the
underlying multigraph of $B$.  Denote this cycle by $Q_C$.

At each object vertex on this cycle, the half-edge belonging to the internal
edge has sign $-$, while the half-edge belonging to the incidence edge has
sign $+$.  At each incidence vertex $p_{v_j,e_j}$, one cycle edge leads to a
vertex pair and has half-edge sign $+$ there, while the other leads to an edge
pair and has half-edge sign $-$ there.  The two relevant half-edge signs are
therefore different at every vertex, so Definition~\ref{def:bidirected} shows
that $Q_C$ is bidirected.  Its associated vertex and edge sets are precisely
those of $C$, and hence $\pi(Q_C)=C$.

Finally, the first part of the proof shows that $Q$ contains $v^0$ if and only
if it contains the whole pair $\{v^0,v^1\}$, which holds if and only if
$\pi(Q)$ contains $v$.  Since $T$ contains precisely one auxiliary vertex
$v^0$ for each $v\in S$, summing this equivalence over $v\in S$ yields
$|T\cap V(Q)|=|S\cap V(\pi(Q))|$.
\end{proof}

\end{document}
